\documentclass{article}
\usepackage{anyfontsize}
\usepackage[UTF8]{ctex} % 如果需要支持中文
\usepackage{fontspec} 

\usepackage[a4paper, left=0.1in, right=0.1in, top=0.05in, bottom=0.05in]{geometry} % 设置四边页边距为0.1英寸{geometry} % 设置页边距为0.5英寸
\usepackage{enumitem} % 用于自定义列表
\usepackage{amsmath}  % 数学公式支持

\setlength{\parindent}{0pt} % 不缩进
\usepackage{etoolbox} % 用于列表操作
%调整类代码
\usepackage{comment}
\usepackage{tocloft} % 用于定制目录样式
\usepackage{hyperref} %不需要目录盒子注释这
\usepackage{ifthen}
\usepackage{tikz}
\usepackage{titletoc}
\usepackage{graphicx}
\usepackage{amssymb}  % 提供数学符号，包括\therefore
%目录设定
%快捷键 CMD + / 注释
%  \renewcommand{\cftdot}{} % 隐藏目录中的页码
%  \cftpagenumbersoff{section}
%  \cftpagenumbersoff{subsection}
%  \makeatletter
%  \renewcommand{\@pnumwidth}{0em}  % 设置页码宽度为0
%  \renewcommand{\@tocrmarg}{0em}   % 设置右边距为0
%  \makeatother
 
%\setCJKmainfont[
 %   Path = C:/USERS/11252/APPDATA/LOCAL/MICROSOFT/WINDOWS/FONTS/,
 %   UprightFont = {SourceHanSerifCN-Light-5.otf},
 %   BoldFont = {SourceHanSerifCN-Medium-6.otf},
%    Scale=1
%]{SourceHanSerifCN}

%\setmainfont{SourceHanSerif}  % 设置英文字体

% 处理冲突的命令
\let\oldbackepsilon\backepsilon
\let\backepsilon\relax
\let\oldeth\eth
\let\eth\relax
\let\olddigamma\digamma
\let\digamma\relax

\usepackage{float}
\usepackage[a4paper, left=0.1in, right=0.1in, top=0.05in, bottom=0.05in]{geometry} % 设置四边页边距为0.1英寸{geometry} % 设置页边距为0.5英寸
\usepackage{enumitem} % 用于自定义列表
\usepackage{amsmath}  % 数学公式支持

\setlength{\parindent}{0pt} % 不缩进
\usepackage{etoolbox} % 用于列表操作
%调整类代码
\usepackage{comment}
\usepackage{tocloft} % 用于定制目录 样式
\usepackage{hyperref} %不需要目录盒子注释这
\usepackage{ifthen}
\usepackage{tikz}
\usepackage{titletoc}
\usepackage{graphicx}
\usepackage{amssymb}  % 提供数学符号，包括\therefore
%目录设定
%快捷键 CMD + / 注释
%  \renewcommand{\cftdot}{} % 隐藏目录中的页码
%  \cftpagenumbersoff{section}
%  \cftpagenumbersoff{subsection}
%  \makeatletter
%  \renewcommand{\@pnumwidth}{0em}  % 设置页码宽度为0
%  \renewcommand{\@tocrmarg}{0em}   % 设置右边距为0
%  \makeatother
 


\usepackage{environ}

\newif\ifshowcontent  % 定义一个条件判断变量
\showcontenttrue    % 设置为 false，隐藏所有 question 环境的内容

\newcounter{question} % 题目计数器
\setcounter{question}{0}
\NewEnviron{question}{%
    \ifshowcontent
        \par\stepcounter{question}\textbf{\thequestion.} \BODY\par % 如果显示内容，则输出
    \fi
}

\NewEnviron{choices}{%
    \ifshowcontent
        \begin{enumerate}[label=\Alph*., leftmargin=*]
            \BODY
        \end{enumerate}\par
    \fi
}

\NewEnviron{correctanswer}{%
    \ifshowcontent
        \textbf{答案:}\BODY\par
    \fi
}



\newenvironment{explanation}
{\textbf{解析:}\par}
{\par}



% 新增考察方面命令
\newcommand{\checklist}[1]{
    \textbf{考察:} #1 \par % 在题目下方显示考察方面
    \addcontentsline{toc}{subsection}{题号\thequestion 考察: #1} % 将考察内容添加到目录
    \vspace{2em} % 添加1em的垂直空间
}

\graphicspath{{images/}} %统一

\begin{document}
 %\fontsize{22}{31}\selectfont
 \tableofcontents % 添加目录
\newpage
%\pagestyle{empty}




\section{考点1:信用风险管理}
\setcounter{question}{0}


\begin{question}
   
\end{question}

\begin{choices}

        \item 

\end{choices}


\begin{correctanswer}
    B
\end{correctanswer}

\begin{explanation}
    \textbf{A选项错误。}
    \begin{itemize}
        \item 
    \end{itemize}
    
    \textbf{B选项正确。}
    \begin{itemize}
        \item 
    \end{itemize}

    \textbf{C选项错误。}
    \begin{itemize}
        \item 
    \end{itemize}
    
    \textbf{D选项错误。}
    \begin{itemize}
        \item 
    \end{itemize}
\end{explanation}

\checklist{}

\newpage


\section{考点2:衍生品信用风险}
\setcounter{question}{0}



\begin{question}
    
\end{question}

\begin{choices}
    \item 
\end{choices}

\begin{correctanswer}
    B
\end{correctanswer}

\begin{explanation}
    \textbf{A选项错误。}
    \begin{itemize}
        \item 
    \end{itemize}
    \textbf{B选项错误。}
    \begin{itemize}
        \item 
    \end{itemize}
    \textbf{C选项错误。}
    \begin{itemize}
        \item 
    \end{itemize}
    \textbf{D选项错误。}
    \begin{itemize}
        \item 
    \end{itemize}
\end{explanation}
\checklist{}




\newpage


\section{考点3:结构化信用产品}
\setcounter{question}{0}
\begin{question}
    Which of the following statements regarding frictions in the securitization of subprime mortgages is correct?
\end{question}

\begin{choices}
    \item The arranger will typically have an information advantage over the originator with regard to the quality of the loans securitized
    \item The originator will typically have an information advantage over the arranger, which can create an incentive for the originator to collaborate with the borrower in filing false loan applications
    \item The major credit rating agencies are paid by investors for their rating service of mortgage-backed securities, and this creates a potential conflict of interest
    \item The use of escrow accounts for insurance and tax payments eliminates the risk of foreclosure
\end{choices}

\begin{correctanswer}
    B
\end{correctanswer}

\begin{explanation}
    \textbf{A选项错误。}
    这个说法不正确：
    \begin{itemize}
        \item 安排人通常不会比发起人更了解贷款质量
        \item 发起人作为贷款的直接发放者具有信息优势
        \item 这个选项误解了信息不对称的方向
    \end{itemize}

    \textbf{B选项正确。}
    这个说法正确：
    \begin{itemize}
        \item 发起人确实比安排人更了解贷款质量
        \item 这种信息优势可能导致发起人与借款人合谋
        \item 发起人可能协助借款人伪造贷款申请材料
        \item 这种欺诈行为会增加证券化产品的风险
    \end{itemize}

    \textbf{C选项错误。}
    这个说法不正确：
    \begin{itemize}
        \item 评级机构不是由投资者支付评级费用
        \item 评级费用通常由证券化产品的发行方支付
        \item 这个选项误解了评级机构的收费模式
    \end{itemize}

    \textbf{D选项错误。}
    这个说法不正确：
    \begin{itemize}
        \item 托管账户只能延缓但不能消除止赎风险
        \item 这个选项夸大了托管账户的作用
        \item 这个选项不符合实际情况
    \end{itemize}
\end{explanation}

\checklist{次贷证券化中发起人具有信息优势，可能引发欺诈风险}


\begin{question}
    Assume there are 100 identical loans with a principal balance of \$500,000 each. Based on a credit analysis, a 300 basis point spread is applied to the borrowers. LIBOR is currently 4\% and the coupon rate will reset annually. The senior, junior, and equity tranches are 75\%, 20\%, and 5\% of the pool, respectively. The spreads on the senior and mezzanine tranches are 2\% and 6\%. Excess cash flow is diverted to trust account above \$1,000,000. Assume the default rate is 2\%. What are the cash flows to the mezzanine and trust account in the first period?
\end{question}

\begin{choices}
    \item \$1,000,000 \$0
    \item \$1,000,000 \$180,000
    \item \$2,250,000 \$200,000
    \item \$2,250,000 \$250,000
\end{choices}

\begin{correctanswer}
    A
\end{correctanswer}

\begin{explanation}
    \textbf{A选项正确。}
    这个说法正确：
    \begin{itemize}
        \item 贷款利率 = 4\%(LIBOR) + 3\%(利差) = 7\%
        \item 第一期总现金流 = 100×\$500,000×7\%×(1-2\%) = \$3,430,000
        \item 高级层现金流 = \$50M×75\%×(4\%+2\%) = \$2,250,000
        \item 夹层级现金流 = \$50M×20\%×(4\%+6\%) = \$1,000,000
        \item 剩余现金流 = \$3,430,000 - \$2,250,000 - \$1,000,000 = \$180,000
        \item 由于剩余现金流小于\$1,000,000，全部归股权层，信托账户为0
    \end{itemize}

    \textbf{B选项错误。}
    这个说法不正确：
    \begin{itemize}
        \item 夹层级现金流计算正确
        \item 但剩余现金流应归股权层，不应进入信托账户
        \item 这个选项误解了现金流分配规则
    \end{itemize}

    \textbf{C选项错误。}
    这个说法不正确：
    \begin{itemize}
        \item 高级层现金流计算正确
        \item 但夹层级现金流计算错误
        \item 这个选项完全误解了现金流分配
    \end{itemize}

    \textbf{D选项错误。}
    这个说法不正确：
    \begin{itemize}
        \item 高级层现金流计算正确
        \item 但夹层级现金流和信托账户金额都计算错误
        \item 这个选项完全误解了现金流分配
    \end{itemize}
\end{explanation}

\checklist{证券化产品现金流分配遵循优先级顺序，剩余现金流归股权层}


\begin{question}
    Assume a firm issues only three capital claims: zero-coupon senior debt with face value of \$300 million; zero-coupon junior debt with face value of \$500 million; and \$200 million in equity. What is the credit enhancement provided to the senior debt?
\end{question}

\begin{choices}
    \item \$200 million
    \item \$300 million
    \item \$500 million
    \item \$700 million
\end{choices}

\begin{correctanswer}
    D
\end{correctanswer}

\begin{explanation}
    \textbf{A选项错误。}
    这个说法不正确：
    \begin{itemize}
        \item 仅考虑了股权价值作为信用增级
        \item 忽略了次级债的信用增级作用
        \item 这个选项低估了信用增级总额
    \end{itemize}

    \textbf{B选项错误。}
    这个说法不正确：
    \begin{itemize}
        \item 仅考虑了高级债的面值
        \item 这个选项完全误解了信用增级的概念
        \item 这个选项与问题无关
    \end{itemize}

    \textbf{C选项错误。}
    这个说法不正确：
    \begin{itemize}
        \item 仅考虑了次级债的面值
        \item 忽略了股权的信用增级作用
        \item 这个选项低估了信用增级总额
    \end{itemize}

    \textbf{D选项正确。}
    这个说法正确：
    \begin{itemize}
        \item 信用增级来自次级债和股权的价值
        \item 次级债面值\$500百万
        \item 股权价值\$200百万
        \item 总信用增级 = \$500百万 + \$200百万 = \$700百万
        \item 这些资金在高级债之前承担损失
    \end{itemize}
\end{explanation}

\checklist{高级债的信用增级来自次级债和股权的价值总和}


\begin{question}
    A pool of high yield bonds is placed in an SPV and three tranches (including the equity tranche) of bonds are issued collateralized by the bonds to create a Collateralized Bond Obligation (CBO). Which of the following is true?
\end{question}

\begin{choices}
    \item At fair value, the value of the issued bonds should be less than the collateral
    \item At fair value, the total default probability, weighted by size of issue, of the issued bonds should equal the default probability of the collateral pool
    \item The equity tranche of the CBO has the least risk of default
    \item The yield on the low risk tranche must be greater than the yield on the collateral pool
\end{choices}

\begin{correctanswer}
    B
\end{correctanswer}

\begin{explanation}
    \textbf{A选项错误。}
    这个说法不正确：
    \begin{itemize}
        \item 在公允价值下，发行的债券价值应等于抵押品价值
        \item 这个选项误解了CBO的基本定价原则
        \item 这个选项不符合无套利原则
    \end{itemize}

    \textbf{B选项正确。}
    这个说法正确：
    \begin{itemize}
        \item CBO和基础证券必须具有相同的市场价值
        \item 加权后的违约概率必须相等
        \item 现金流模式必须相似
        \item 风险特征必须一致
    \end{itemize}

    \textbf{C选项错误。}
    这个说法不正确：
    \begin{itemize}
        \item 股权层级具有最高的违约风险
        \item 这个选项完全误解了CBO的风险分配
        \item 这个选项与实际情况相反
    \end{itemize}

    \textbf{D选项错误。}
    这个说法不正确：
    \begin{itemize}
        \item 低风险层级的收益率必须低于抵押品池的收益率
        \item 这个选项误解了CBO的收益分配
        \item 这个选项与实际情况相反
    \end{itemize}
\end{explanation}

\checklist{CBO与基础证券必须具有相同的市场价值和风险特征}


\begin{question}
    Continuously increasing default probability (while holding default correlation constant) will most likely have what effect on the credit VaR of mezzanine and equity tranches?
\end{question}

\begin{choices}
    \item Increase Increase then decrease
    \item Increase Decrease then increase
    \item Decrease Increase then decrease
    \item Decrease Decrease then increase
\end{choices}

\begin{correctanswer}
    C
\end{correctanswer}

\begin{explanation}
    \textbf{A选项错误。}
    这个说法不正确：
    \begin{itemize}
        \item 股权层VaR会随着违约概率上升而下降
        \item 夹层级VaR的变化趋势描述错误
        \item 这个选项完全误解了VaR的变化特征
    \end{itemize}

    \textbf{B选项错误。}
    这个说法不正确：
    \begin{itemize}
        \item 股权层VaR会随着违约概率上升而下降
        \item 夹层级VaR的变化趋势描述错误
        \item 这个选项完全误解了VaR的变化特征
    \end{itemize}

    \textbf{C选项正确。}
    这个说法正确：
    \begin{itemize}
        \item 股权层VaR下降原因：
        \begin{itemize}
            \item 违约概率上升时，股权层损失确定性增加
            \item 损失波动性减小，导致VaR下降
        \end{itemize}
        \item 夹层级VaR先升后降原因：
        \begin{itemize}
            \item 低违约概率时，表现类似高级层，VaR上升
            \item 高违约概率时，表现类似股权层，VaR下降
        \end{itemize}
    \end{itemize}

    \textbf{D选项错误。}
    这个说法不正确：
    \begin{itemize}
        \item 股权层VaR变化趋势正确
        \item 但夹层级VaR的变化趋势描述错误
        \item 这个选项部分正确但整体错误
    \end{itemize}
\end{explanation}

\checklist{违约概率上升时，股权层VaR下降，夹层级VaR先升后降}


\begin{question}
    Which of the following statements on credit enhancement on securitized assets for the purposes of mitigating credit risk and liquidation risk are true?
    I. Internal and external enhancements can be used for both credit and liquidity risk.
    II. Liquidity reserves for liquidity risk are similar to a cash collateral account for credit risk.
\end{question}

\begin{choices}
    \item I only
    \item II only
    \item Both I and II
    \item Neither I nor II
\end{choices}

\begin{correctanswer}
    C
\end{correctanswer}

\begin{explanation}
    \textbf{A选项错误。}
    这个说法不正确：
    \begin{itemize}
        \item 仅考虑了内部和外部增级措施
        \item 忽略了流动性储备与现金抵押账户的相似性
        \item 这个选项不完整
    \end{itemize}

    \textbf{B选项错误。}
    这个说法不正确：
    \begin{itemize}
        \item 仅考虑了流动性储备与现金抵押账户的相似性
        \item 忽略了内部和外部增级措施
        \item 这个选项不完整
    \end{itemize}

    \textbf{C选项正确。}
    这个说法正确：
    \begin{itemize}
        \item 陈述I正确：
        \begin{itemize}
            \item 信用风险和流动性风险都可以使用内部增级
            \item 信用风险和流动性风险都可以使用外部增级
        \end{itemize}
        \item 陈述II正确：
        \begin{itemize}
            \item 流动性储备用于应对流动性风险
            \item 现金抵押账户用于应对信用风险
            \item 两者在功能上具有相似性
        \end{itemize}
    \end{itemize}

    \textbf{D选项错误。}
    这个说法不正确：
    \begin{itemize}
        \item 两个陈述都是正确的
        \item 这个选项完全错误
        \item 这个选项误解了增级措施的特征
    \end{itemize}
\end{explanation}

\checklist{证券化产品可通过内部外部增级措施应对信用和流动性风险}


\begin{question}
    An important step in structuring a securitization is determining an adequate level of credit support or enhancement. Which of the following is not one of common credit enhancement forms for securitization transactions?
\end{question}

\begin{choices}
    \item Subordinated tranches of securitization debt
    \item Excess Spread account
    \item Cash reserve account
    \item Sinking fund account
\end{choices}

\begin{correctanswer}
    D
\end{correctanswer}

\begin{explanation}
    \textbf{A选项错误。}
    这个说法不正确：
    \begin{itemize}
        \item 次级层级是常见的信用增级形式
        \item 通过优先级顺序提供信用保护
        \item 这个选项是有效的信用增级方式
    \end{itemize}

    \textbf{B选项错误。}
    这个说法不正确：
    \begin{itemize}
        \item 超额利差账户是常见的信用增级形式
        \item 通过资产池收益与支出差额提供保护
        \item 这个选项是有效的信用增级方式
    \end{itemize}

    \textbf{C选项错误。}
    这个说法不正确：
    \begin{itemize}
        \item 现金储备账户是常见的信用增级形式
        \item 通过预留现金提供信用保护
        \item 这个选项是有效的信用增级方式
    \end{itemize}

    \textbf{D选项正确。}
    这个说法正确：
    \begin{itemize}
        \item 偿债基金不是信用增级形式
        \item 偿债基金用于定期赎回债券
        \item 偿债基金是债券发行的常规安排
        \item 偿债基金与信用增级无关
    \end{itemize}
\end{explanation}

\checklist{偿债基金不是证券化产品的信用增级形式}

\begin{question}
    Which of the following is not a friction in the subprime securitization market?
\end{question}

\begin{choices}
    \item Investor and rating agency
    \item Servicer and mortgagor
    \item Mortgagor and arranger
    \item Asset manager and investor
\end{choices}

\begin{correctanswer}
    C
\end{correctanswer}

\begin{explanation}
    \textbf{A选项错误。}
    这个说法不正确：
    \begin{itemize}
        \item 投资者与评级机构之间存在摩擦
        \item 评级机构可能低估风险
        \item 投资者可能过度依赖评级
        \item 这个选项是次贷证券化市场的摩擦之一
    \end{itemize}

    \textbf{B选项错误。}
    这个说法不正确：
    \begin{itemize}
        \item 服务商与借款人之间存在摩擦
        \item 服务商可能缺乏激励去管理贷款
        \item 借款人可能面临服务问题
        \item 这个选项是次贷证券化市场的摩擦之一
    \end{itemize}

    \textbf{C选项正确。}
    这个说法正确：
    \begin{itemize}
        \item 借款人与承销人没有直接接触
        \item 两者之间不存在直接关系
        \item 因此不存在摩擦
        \item 这个选项不是次贷证券化市场的摩擦
    \end{itemize}

    \textbf{D选项错误。}
    这个说法不正确：
    \begin{itemize}
        \item 资产管理人与投资者之间存在摩擦
        \item 可能存在利益冲突
        \item 可能存在信息不对称
        \item 这个选项是次贷证券化市场的摩擦之一
    \end{itemize}
\end{explanation}

\checklist{次贷证券化市场中借款人与承销人无直接接触，不存在摩擦}


\begin{question}
    A hedge fund is considering taking positions in various tranches of a collateralized debt obligation (CDO). The fund's chief economist predicts that the default probability will decrease significantly and that the default correlation will increase. Based on this prediction, which of the following is a good strategy to pursue?
\end{question}

\begin{choices}
    \item Buy the senior tranche and buy the equity tranche
    \item Buy the senior tranche and sell the equity tranche
    \item Sell the senior tranche and sell equity tranche
    \item Sell the senior tranche and buy the equity tranche
\end{choices}

\begin{correctanswer}
    D
\end{correctanswer}

\begin{explanation}
    \textbf{A选项错误。}
    这个说法不正确：
    \begin{itemize}
        \item 违约概率下降对股权层有利
        \item 但违约相关性上升对高级层不利
        \item 这个策略没有考虑违约相关性的影响
    \end{itemize}

    \textbf{B选项错误。}
    这个说法不正确：
    \begin{itemize}
        \item 违约概率下降对股权层有利
        \item 不应该做空股权层
        \item 这个策略完全误解了市场预期
    \end{itemize}

    \textbf{C选项错误。}
    这个说法不正确：
    \begin{itemize}
        \item 违约概率下降对股权层有利
        \item 不应该做空股权层
        \item 这个策略完全误解了市场预期
    \end{itemize}

    \textbf{D选项正确。}
    这个说法正确：
    \begin{itemize}
        \item 违约概率下降的影响：
        \begin{itemize}
            \item 提高股权层价值
            \item 降低违约风险
        \end{itemize}
        \item 违约相关性上升的影响：
        \begin{itemize}
            \item 提高股权层价值
            \item 降低高级层价值
        \end{itemize}
        \item 因此应该做多股权层，做空高级层
    \end{itemize}
\end{explanation}

\checklist{违约概率下降和相关性上升时，应做多股权层做空高级层}


\begin{question}
    Which of the following subprime characteristics provide direct protection for senior tranches?
\end{question}

\begin{choices}
    \item Subordination, excess spread, and shifting interest
    \item Subordination, prepayments, and shifting interest
    \item Overcollateralization, excess spread, and tinning of losses
    \item Overcollateralization, excess spread, and prepayments
\end{choices}

\begin{correctanswer}
    A
\end{correctanswer}

\begin{explanation}
    \textbf{A选项正确。}
    这个说法正确：
    \begin{itemize}
        \item 分层（Subordination）：
        \begin{itemize}
            \item 通过优先级顺序提供保护
            \item 次级层级先承担损失
        \end{itemize}
        \item 超额利差（Excess Spread）：
        \begin{itemize}
            \item 资产池收益大于支出
            \item 提供额外的信用保护
        \end{itemize}
        \item 利率锁定期（Shifting Interest）：
        \begin{itemize}
            \item 锁定期内本金优先偿付高级层
            \item 中间层只能获得利息
        \end{itemize}
    \end{itemize}

    \textbf{B选项错误。}
    这个说法不正确：
    \begin{itemize}
        \item 提前还款不是直接保护机制
        \item 可能加速或延缓现金流
        \item 这个选项误解了保护机制
    \end{itemize}

    \textbf{C选项错误。}
    这个说法不正确：
    \begin{itemize}
        \item 损失时间不是直接保护机制
        \item 只影响超额利差
        \item 这个选项误解了保护机制
    \end{itemize}

    \textbf{D选项错误。}
    这个说法不正确：
    \begin{itemize}
        \item 提前还款不是直接保护机制
        \item 可能加速或延缓现金流
        \item 这个选项误解了保护机制
    \end{itemize}
\end{explanation}

\checklist{次贷证券化中分层、超额利差和利率锁定期为高级层提供直接保护}


\begin{question}
    Which of the following statements most accurately describes the effect of selling a loan without recourse?
\end{question}

\begin{choices}
    \item The bank that sells the loan retains a contingent liability
    \item The bank that sells the loan bears a specified percentage of the credit risk
    \item The loan is removed from the balance sheet of the bank that sells the loan
    \item The purchaser has the right to sell the loan back to the bank that originated the loan
\end{choices}

\begin{correctanswer}
    C
\end{correctanswer}

\begin{explanation}
    \textbf{A选项错误。}
    这个说法不正确：
    \begin{itemize}
        \item 无追索权出售意味着银行不保留或有负债
        \item 所有信用风险都转移给买方
        \item 这个选项误解了无追索权的含义
    \end{itemize}

    \textbf{B选项错误。}
    这个说法不正确：
    \begin{itemize}
        \item 无追索权出售意味着银行不承担任何信用风险
        \item 所有信用风险都转移给买方
        \item 这个选项误解了无追索权的含义
    \end{itemize}

    \textbf{C选项正确。}
    这个说法正确：
    \begin{itemize}
        \item 无追索权出售的影响：
        \begin{itemize}
            \item 贷款从银行资产负债表中移除
            \item 所有信用风险转移给买方
            \item 银行不再承担任何风险
        \end{itemize}
    \end{itemize}

    \textbf{D选项错误。}
    这个说法不正确：
    \begin{itemize}
        \item 无追索权出售意味着买方不能将贷款卖回给银行
        \item 这个选项误解了无追索权的含义
        \item 这个选项与实际情况相反
    \end{itemize}
\end{explanation}

\checklist{无追索权出售贷款时，贷款从银行资产负债表中移除}


\begin{question}
    During the recent credit crisis, subprime mortgages received pressure from upward movements in interest rates, along with the impossibility for most of these borrowers to refinance. As a result, many subprime borrowers allowed their reduced value homes to be taken over by the lender. All of the following statements describe the effect subprime actions had on mortgage-backed security portfolios except:
\end{question}

\begin{choices}
    \item Repayment issues in lower tranches impacted confidence in higher tranches
    \item There was a fire sale of securitized debt
    \item Panic among MBS investors led to a flight to safer assets
    \item The defaults in the lower tranches were high but did not affect senior tranche investors
\end{choices}

\begin{correctanswer}
    D
\end{correctanswer}

\begin{explanation}
    \textbf{A选项错误。}
    这个说法不正确：
    \begin{itemize}
        \item 次级层级还款问题确实影响了高级层信心
        \item 这是次贷危机的典型特征
        \item 这个选项描述了实际发生的情况
    \end{itemize}

    \textbf{B选项错误。}
    这个说法不正确：
    \begin{itemize}
        \item 确实发生了证券化债务的甩卖
        \item 这是市场恐慌的表现
        \item 这个选项描述了实际发生的情况
    \end{itemize}

    \textbf{C选项错误。}
    这个说法不正确：
    \begin{itemize}
        \item MBS投资者确实出现恐慌
        \item 确实转向更安全的资产
        \item 这个选项描述了实际发生的情况
    \end{itemize}

    \textbf{D选项正确。}
    这个说法正确：
    \begin{itemize}
        \item 这个说法不符合实际情况：
        \begin{itemize}
            \item 次级层级违约确实影响了高级层投资者
            \item 信心危机蔓延到所有层级
            \item 高级层投资者也受到损失
        \end{itemize}
    \end{itemize}
\end{explanation}

\checklist{次贷危机中次级层级违约影响了所有层级的投资者}


\begin{question}
    Which of the following statements about portfolio losses and default correlation are most likely correct?
    I. Increasing default correlation decreases senior tranche values but increases equity tranche values.
    II. At high default rates, increasing default correlation decreases mezzanine bond prices.
\end{question}

\begin{choices}
    \item I only
    \item II only
    \item Both I and II
    \item Neither I nor II
\end{choices}

\begin{correctanswer}
    A
\end{correctanswer}

\begin{explanation}
    \textbf{A选项正确。}
    这个说法正确：
    \begin{itemize}
        \item 陈述I正确：
        \begin{itemize}
            \item 违约相关性上升增加极端组合收益的可能性
            \item 高违约可能性降低高级层价值
            \item 低违约可能性提高股权层价值
        \end{itemize}
        \item 陈述II错误：
        \begin{itemize}
            \item 高违约率时，相关性上升增加极端收益可能性
            \item 这对股权层和夹层级都有利
            \item 不会降低夹层级债券价格
        \end{itemize}
    \end{itemize}

    \textbf{B选项错误。}
    这个说法不正确：
    \begin{itemize}
        \item 陈述I是正确的
        \item 陈述II是错误的
        \item 这个选项部分正确但整体错误
    \end{itemize}

    \textbf{C选项错误。}
    这个说法不正确：
    \begin{itemize}
        \item 陈述I是正确的
        \item 陈述II是错误的
        \item 这个选项部分正确但整体错误
    \end{itemize}

    \textbf{D选项错误。}
    这个说法不正确：
    \begin{itemize}
        \item 陈述I是正确的
        \item 陈述II是错误的
        \item 这个选项完全错误
    \end{itemize}
\end{explanation}

\checklist{违约相关性上升降低高级层价值但提高股权层价值}


\begin{question}
    How many of the following statements concerning the capital structure in a securitization are most likely correct?
    I. The mezzanine tranche is typically the smallest tranche size.
    II. The mezzanine and equity tranches typically offer fixed coupons.
    III. The senior tranche typically receives the lowest coupon.
\end{question}

\begin{choices}
    \item No statements are correct
    \item One statement is correct
    \item Two statements are correct
    \item Three statements are correct
\end{choices}

\begin{correctanswer}
    B
\end{correctanswer}

\begin{explanation}

    \textbf{B选项正确。}
    这个说法正确：
    \begin{itemize}
        \item 陈述I错误：
        \begin{itemize}
            \item 夹层级不是最小的层级
            \item 股权层通常是最小的层级
        \end{itemize}
        \item 陈述II错误：
        \begin{itemize}
            \item 股权层没有固定票息
            \item 股权层获得剩余现金流
        \end{itemize}
        \item 陈述III正确：
        \begin{itemize}
            \item 高级层被认为最安全
            \item 因此获得最低票息
        \end{itemize}
    \end{itemize}

\end{explanation}

\checklist{证券化产品中高级层获得最低票息，股权层最小且无固定票息}


\begin{question}
    A firm is experiencing financial difficulties. Using a contingent claims approach, which of the following best describes the valuation of their senior and subordinated debt?
\end{question}

\begin{choices}
    \item Both the senior debt and subordinated debt have positive exposures to debt maturity, firm volatility, and interest rates
    \item The senior debt has negative exposures to debt maturity, firm volatility, and interest rates. The subordinated debt has positive exposures to debt maturity, firm volatility, and interest rates
    \item The senior debt has positive exposures to debt maturity, firm volatility, and interest rates. The subordinated debt has negative exposures to debt maturity, firm volatility, and interest rates
    \item Both the senior debt and subordinated debt have negative exposures to debt maturity, firm volatility, and interest rates
\end{choices}

\begin{correctanswer}
    B
\end{correctanswer}

\begin{explanation}
    \textbf{A选项错误。}
    这个说法不正确：
    \begin{itemize}
        \item 高级债对这些因素有负向敞口
        \item 这个选项误解了高级债的特征
        \item 这个选项完全错误
    \end{itemize}

    \textbf{B选项正确。}
    这个说法正确：
    \begin{itemize}
        \item 高级债特征：
        \begin{itemize}
            \item 对到期期限有负向敞口
            \item 对公司波动率有负向敞口
            \item 对利率有负向敞口
        \end{itemize}
        \item 次级债特征：
        \begin{itemize}
            \item 在财务困境时表现类似股权
            \item 对到期期限有正向敞口
            \item 对公司波动率有正向敞口
            \item 对利率有正向敞口
        \end{itemize}
    \end{itemize}

    \textbf{C选项错误。}
    这个说法不正确：
    \begin{itemize}
        \item 高级债对这些因素有负向敞口
        \item 次级债在财务困境时对这些因素有正向敞口
        \item 这个选项完全误解了债务特征
    \end{itemize}

    \textbf{D选项错误。}
    这个说法不正确：
    \begin{itemize}
        \item 次级债在财务困境时对这些因素有正向敞口
        \item 这个选项误解了次级债的特征
        \item 这个选项部分正确但整体错误
    \end{itemize}
\end{explanation}

\checklist{财务困境时高级债有负向敞口，次级债有正向敞口}


\begin{question}
    Which of the following methods is typically implemented as the second stage of principal component analysis?
\end{question}

\begin{choices}
    \item Factor analysis
    \item Divisive clustering
    \item Hierarchical clustering
    \item The canonical correlation method
\end{choices}

\begin{correctanswer}
    A
\end{correctanswer}

\begin{explanation}
    \textbf{A选项正确。}
    这个说法正确：
    \begin{itemize}
        \item 主成分分析的第二阶段：
        \begin{itemize}
            \item 第一步是对主成分进行标准化
            \item 第二步是对新变量（因子载荷）进行标准化
            \item 因子分析是常用的第二阶段方法
        \end{itemize}
    \end{itemize}

    \textbf{B选项错误。}
    这个说法不正确：
    \begin{itemize}
        \item 分裂聚类不是主成分分析的第二阶段
        \item 这是另一种聚类方法
        \item 这个选项与主成分分析无关
    \end{itemize}

    \textbf{C选项错误。}
    这个说法不正确：
    \begin{itemize}
        \item 层次聚类不是主成分分析的第二阶段
        \item 这是另一种聚类方法
        \item 这个选项与主成分分析无关
    \end{itemize}

    \textbf{D选项错误。}
    这个说法不正确：
    \begin{itemize}
        \item 典型相关方法不是主成分分析的第二阶段
        \item 这是另一种统计分析方法
        \item 这个选项与主成分分析无关
    \end{itemize}
\end{explanation}

\checklist{主成分分析的第二阶段通常使用因子分析方法}


\begin{question}
    Under the contingent claim approach to the firm's capital structure, which of the following statements is true? Assume the amount of senior debt, subordinated debt, and equity is represented as F, U, and S, respectively.
\end{question}

\begin{choices}
    \item The value of subordinated debt is less than the value of senior debt
    \item Subordinated debt can be represented by a long call with exercise price of F and short call
    \item Subordinated debt behaves more like equity in distress and more like debt when the firm is not in distress
    \item The value of subordinated debt is always greater than the value of equity
\end{choices}

\begin{correctanswer}
    C
\end{correctanswer}

\begin{explanation}
    \textbf{A选项错误。}
    这个说法不正确：
    \begin{itemize}
        \item 次级债价值不一定小于高级债
        \item 相对金额可能差异很大
        \item 这个选项过于绝对
    \end{itemize}

    \textbf{B选项错误。}
    这个说法不正确：
    \begin{itemize}
        \item 次级债的期权表示错误
        \item 应该是执行价格为F的看涨期权多头
        \item 和执行价格为F+U的看涨期权空头
        \item 这个选项误解了期权组合
    \end{itemize}

    \textbf{C选项正确。}
    这个说法正确：
    \begin{itemize}
        \item 财务困境时：
        \begin{itemize}
            \item 股权价值较小
            \item 次级债表现类似股权
            \item 更接近获得剩余现金流
        \end{itemize}
        \item 非财务困境时：
        \begin{itemize}
            \item 公司价值较高
            \item 次级债表现类似传统债务
        \end{itemize}
    \end{itemize}

    \textbf{D选项错误。}
    这个说法不正确：
    \begin{itemize}
        \item 次级债价值不一定大于股权
        \item 相对金额可能差异很大
        \item 这个选项过于绝对
    \end{itemize}
\end{explanation}

\checklist{财务困境时次级债表现类似股权，非困境时表现类似债务}


\begin{question}
    Harrison Michaels, FRM, an analyst at Hudson Risk Analytics, is discussing the default sensitivities of equity, mezzanine, and senior tranches and makes the following statements. Which of the following statements is(are) most likely correct?
    I. Default sensitivities are largest close to the attachment point between tranches.
    II. Default sensitivities are computed by shocking the credit default swap (CDS) default curve.
\end{question}

\begin{choices}
    \item I only
    \item II only
    \item Both I and II
    \item Neither I nor II
\end{choices}

\begin{correctanswer}
    A
\end{correctanswer}

\begin{explanation}
    \textbf{A选项正确。}
    这个说法正确：
    \begin{itemize}
        \item 陈述I正确：
        \begin{itemize}
            \item 违约敏感性在层级连接点附近最大
            \item 这是因为损失接近连接点时影响最大
            \item 所有层级的违约敏感性都是正值
            \item 高违约率时敏感性趋近于零
        \end{itemize}
        \item 陈述II错误：
        \begin{itemize}
            \item 违约敏感性不是通过冲击CDS违约曲线计算
            \item 而是通过冲击各种假设的违约概率计算
            \item 这个陈述误解了计算方法
        \end{itemize}
    \end{itemize}

\end{explanation}

\checklist{违约敏感性在层级连接点附近最大，通过冲击违约概率计算}

\begin{question}
    Advanced Pharmaceuticals is considering an investment in a very risky drug therapy treatment. If the investment is successful, Advanced Pharmaceuticals can earn substantial profits. On the other hand, if the therapy treatments create long-term side-effects, the firm will be subject to expensive litigation. How should Advanced Pharmaceuticals structure its investment to minimize its cost of capital?
\end{question}

\begin{choices}
    \item Ring-fence the investment so a special purpose entity (SPE) can issue debt at lower costs
    \item Ring-fence the investment into a SPE so the parent company can increase transparency
    \item Securitize the investment due to its predictable cash flows
    \item Sell the investment in a true sale, increasing adverse selection
\end{choices}

\begin{correctanswer}
    B
\end{correctanswer}

\begin{explanation}
    \textbf{A选项错误。}
    这个说法不正确：
    \begin{itemize}
        \item SPE本身包含高风险投资
        \item 发行债务成本会更高
        \item 这个选项误解了SPE的作用
    \end{itemize}

    \textbf{B选项正确。}
    这个说法正确：
    \begin{itemize}
        \item 通过SPE隔离风险投资：
        \begin{itemize}
            \item 母公司剩余资产更透明
            \item 可以以更优惠条件发行债务
            \item 降低资本成本
        \end{itemize}
        \item 真实销售可以：
        \begin{itemize}
            \item 在法律上分离风险投资
            \item 提高透明度
            \item 减少逆向选择
        \end{itemize}
    \end{itemize}

    \textbf{C选项错误。}
    这个说法不正确：
    \begin{itemize}
        \item 该投资没有可预测的现金流
        \item 不像抵押贷款或汽车贷款
        \item 这个选项误解了证券化的条件
    \end{itemize}

    \textbf{D选项错误。}
    这个说法不正确：
    \begin{itemize}
        \item 真实销售会减少逆向选择
        \item 不会增加逆向选择
        \item 这个选项误解了真实销售的影响
    \end{itemize}
\end{explanation}

\checklist{通过SPE隔离风险投资可提高透明度，降低资本成本}


\begin{question}
    Harris Smith, CFO of XYZ Bank Corp, is considering a \$500 million loan securitization. He has enlisted a well-respected structuring agent to help decide on the most beneficial structure. XYZ is a \$100 billion regional bank with a moderately strong balance sheet. Its current credit rating on unsecured debt is BBB. It recently issued a secured bond issue with a credit rating of A after ring-fencing certain assets. XYZ desires to minimize the cost of funds and achieve AAA credit rating on the senior tranche of the new securitization. After reviewing the financials of XYZ and forecasting future economic conditions, the structure has recommended an arbitrage CDO with the following loss distributions: Equity tranche: 0-30\%. Junior tranche: 30-50\%. Smith should use which of the following CDO structures?
\end{question}

\begin{choices}
    \item Arbitrage CDO with \$25 million equity tranche
    \item Arbitrage CDO with \$150 million equity tranche
    \item Balance sheet CDO with \$25 million equity tranche
    \item Balance sheet CDO with \$150 million equity tranche
\end{choices}

\begin{correctanswer}
    C
\end{correctanswer}

\begin{explanation}
    \textbf{A选项错误。}
    这个说法不正确：
    \begin{itemize}
        \item 银行要证券化自己发放的贷款
        \item 应该使用资产负债表CDO
        \item 这个选项选择了错误的CDO类型
    \end{itemize}

    \textbf{B选项错误。}
    这个说法不正确：
    \begin{itemize}
        \item 银行要证券化自己发放的贷款
        \item 应该使用资产负债表CDO
        \item 股权层级过大
        \item 这个选项完全错误
    \end{itemize}

    \textbf{C选项正确。}
    这个说法正确：
    \begin{itemize}
        \item 资产负债表CDO：
        \begin{itemize}
            \item 适合证券化银行自身贷款
            \item 符合银行的需求
        \end{itemize}
        \item 股权层级：
        \begin{itemize}
            \item \$25百万占发行总额的5\%
            \item 规模适中，可以保持高级层AAA评级
            \item 有助于最小化融资成本
        \end{itemize}
    \end{itemize}

    \textbf{D选项错误。}
    这个说法不正确：
    \begin{itemize}
        \item 资产负债表CDO类型正确
        \item 但股权层级过大
        \item 可能影响高级层评级
        \item 这个选项部分正确但整体错误
    \end{itemize}
\end{explanation}

\checklist{银行应使用资产负债表CDO，股权层级占5\%以保持高级层AAA评级}


\begin{question}
    Which of the following statements regarding subprime mortgages is true?
    I. Senior tranches receive the highest return due to overcollateralization.
    II. Corporate bonds and subprime pools with the same credit rating will exhibit the same variation in losses.
\end{question}

\begin{choices}
    \item I only
    \item II only
    \item Both I and II
    \item Neither I nor II
\end{choices}

\begin{correctanswer}
    D
\end{correctanswer}

\begin{explanation}
    \textbf{A选项错误。}
    这个说法不正确：
    \begin{itemize}
        \item 陈述I错误：
        \begin{itemize}
            \item 高级层获得最低收益
            \item 这是因为其在资本结构中的优先级
            \item 超额抵押提供保护而非提高收益
        \end{itemize}
        \item 陈述II错误：
        \begin{itemize}
            \item 相同评级下损失波动性不同
            \item ABS结构损失波动性更大
        \end{itemize}
    \end{itemize}

    \textbf{B选项错误。}
    这个说法不正确：
    \begin{itemize}
        \item 陈述I错误
        \item 陈述II错误
        \item 这个选项完全错误
    \end{itemize}

    \textbf{C选项错误。}
    这个说法不正确：
    \begin{itemize}
        \item 陈述I错误
        \item 陈述II错误
        \item 这个选项完全错误
    \end{itemize}

    \textbf{D选项正确。}
    这个说法正确：
    \begin{itemize}
        \item 陈述I错误：
        \begin{itemize}
            \item 高级层获得最低收益
            \item 超额抵押提供保护而非提高收益
        \end{itemize}
        \item 陈述II错误：
        \begin{itemize}
            \item 相同评级下损失波动性不同
            \item ABS结构损失波动性更大
        \end{itemize}
    \end{itemize}
\end{explanation}

\checklist{次贷证券化中高级层收益最低，ABS结构损失波动性更大}



\begin{question}
    The manager makes additional observations as follows:
    • The loans in the collateral pool pay a fixed spread of 2.2% over the swap curve
    • There were no defaults in year 1
    • The value of the overcollateralization account at the end of year 1 was EUR 0
    What is the value of the overcollateralization account at the end of year 2 if there are 8 defaults in year 2?
\end{question}

\begin{choices}
    \item EUR 600,000
    \item EUR 1,056,000
    \item EUR 2,544,000
    \item EUR 3,600,000
\end{choices}

\begin{correctanswer}
    C
\end{correctanswer}

\begin{explanation}
    \textbf{A选项错误。}
    这个说法不正确：
    \begin{itemize}
        \item 计算错误
        \item 低估了超额抵押账户的价值
        \item 这个选项完全错误
    \end{itemize}

    \textbf{B选项错误。}
    这个说法不正确：
    \begin{itemize}
        \item 计算错误
        \item 低估了超额抵押账户的价值
        \item 这个选项完全错误
    \end{itemize}

    \textbf{C选项正确。}
    这个说法正确：
    \begin{itemize}
        \item 计算过程：
        \begin{itemize}
            \item 贷款利差为2.2%
            \item 第一年无违约
            \item 第二年8笔违约
            \item 超额抵押账户价值 = EUR 2,544,000
        \end{itemize}
    \end{itemize}

    \textbf{D选项错误。}
    这个说法不正确：
    \begin{itemize}
        \item 计算错误
        \item 高估了超额抵押账户的价值
        \item 这个选项完全错误
    \end{itemize}
\end{explanation}

\checklist{超额抵押账户价值取决于贷款利差和违约情况}




\newpage



\section{考点4:证券化}
\setcounter{question}{0}

\begin{question}
    Assume an MBS is composed of the following four different pools of mortgages: \$2 million of mortgages that have a maturity of 90 days. \$3 million of mortgages that have a maturity of 180 days. \$5 million of mortgages that have a maturity of 270 days, \$10 million of mortgages that have a maturity of 360 days. What is the weighted average maturity (WAM) of these mortgage pools?
\end{question}

\begin{choices}
    \item 167 days
    \item 225 days
    \item 252 days
    \item 284 days
\end{choices}

\begin{correctanswer}
    D
\end{correctanswer}

\begin{explanation}
    \textbf{A选项错误。}
    这个说法不正确：
    \begin{itemize}
        \item 计算错误
        \item 低估了加权平均期限
        \item 这个选项完全错误
    \end{itemize}

    \textbf{B选项错误。}
    这个说法不正确：
    \begin{itemize}
        \item 计算错误
        \item 低估了加权平均期限
        \item 这个选项完全错误
    \end{itemize}

    \textbf{C选项错误。}
    这个说法不正确：
    \begin{itemize}
        \item 计算错误
        \item 低估了加权平均期限
        \item 这个选项完全错误
    \end{itemize}

    \textbf{D选项正确。}
    这个说法正确：
    \begin{itemize}
        \item 计算过程：
        \begin{itemize}
            \item 总金额 = \$2M + \$3M + \$5M + \$10M = \$20M
            \item WAM = (90×2 + 180×3 + 270×5 + 360×10) / 20
            \item = (180 + 540 + 1350 + 3600) / 20
            \item = 284天
        \end{itemize}
    \end{itemize}
\end{explanation}

\checklist{抵押贷款支持证券的加权平均期限}


\begin{question}
    Suppose the annual prepayment rate CPR for a mortgage-backed security is 6 percent. What is the corresponding single-monthly mortality rate SMM?
\end{question}

\begin{choices}
    \item 0.514\%
    \item 0.334\%
    \item 0.500\%
    \item 1.355\%
\end{choices}

\begin{correctanswer}
    A
\end{correctanswer}

\begin{explanation}
    \textbf{A选项正确。}
    这个说法正确：
    \begin{itemize}
        \item 计算过程：
             CPR = 6\%
             $SMM = 1 - (1 - CPR)^{1/12}$
             $= 1 - (1 - 0.06)^{1/12}$
             $= 0.514\%$
        
    \end{itemize}


\end{explanation}

\checklist{年提前还款率6\%对应的月提前还款率为0.514\%}

\begin{question}
    Securitized products are often customized to meet the needs of the investor as well as the originator. What type of asset-backed securities (ABSs) typically uses a revolving structure?
\end{question}

\begin{choices}
    \item Residential mortgage
    \item Credit card debt
    \item Commercial mortgage
    \item Commercial paper
\end{choices}

\begin{correctanswer}
    B
\end{correctanswer}

\begin{explanation}
    \textbf{A选项错误。}
    这个说法不正确：
    \begin{itemize}
        \item 住宅抵押贷款不使用循环结构
        \item 有固定的还款计划
        \item 这个选项与循环结构无关
    \end{itemize}

    \textbf{B选项正确。}
    这个说法正确：
    \begin{itemize}
        \item 信用卡债务使用循环结构：
        \begin{itemize}
            \item 没有预先设定的还款计划
            \item 本金以大额方式偿还
            \item 分为循环期和摊销期
            \item 循环期将本金用于购买新贷款
            \item 摊销期进行最终偿付
        \end{itemize}
    \end{itemize}

    \textbf{C选项错误。}
    这个说法不正确：
    \begin{itemize}
        \item 商业抵押贷款不使用循环结构
        \item 有固定的还款计划
        \item 这个选项与循环结构无关
    \end{itemize}

    \textbf{D选项错误。}
    这个说法不正确：
    \begin{itemize}
        \item 商业票据不使用循环结构
        \item 有固定的到期日
        \item 这个选项与循环结构无关
    \end{itemize}
\end{explanation}

\checklist{信用卡债务证券化通常使用循环结构，分为循环期和摊销期}



\begin{question}
    Which of the following statements regarding credit enhancements in the process of structuring a securitization through a special purpose vehicle (SPV) is correct?
\end{question}

\begin{choices}
    \item The securitization process is structured such that the asset side of the SPV has a lower cost than the liability side of the SPV
    \item Credit enhancements are typically only associated with mortgage-backed securities (MBS) and are not used in other types of asset-backed securities (ABS)
    \item The most senior class of notes is often overcollateralized in order to reduce the risk of the asset-backed security (ABS)
    \item A margin step-up is sometimes used by an asset-backed securities (ABS) where the coupon structure increases after a call date
\end{choices}

\begin{correctanswer}
    D
\end{correctanswer}

\begin{explanation}
    \textbf{A选项错误。}
    这个说法不正确：
    \begin{itemize}
        \item SPV的负债成本低于资产成本
        \item 这样才能产生超额利差
        \item 这个选项与实际情况相反
    \end{itemize}

    \textbf{B选项错误。}
    这个说法不正确：
    \begin{itemize}
        \item 信用增级不仅用于MBS
        \item 也用于其他类型的ABS
        \item 这个选项过于局限
    \end{itemize}

    \textbf{C选项错误。}
    这个说法不正确：
    \begin{itemize}
        \item 最低层级票据通常被超额担保
        \item 不是最高层级
        \item 这个选项误解了超额担保的对象
    \end{itemize}

    \textbf{D选项正确。}
    这个说法正确：
    \begin{itemize}
        \item ABS可能使用利差递增：
        \begin{itemize}
            \item 在赎回日期后提高票息结构
            \item 这是一种信用增级方式
            \item 有助于提高证券的吸引力
        \end{itemize}
    \end{itemize}
\end{explanation}

\checklist{ABS可能使用利差递增作为信用增级方式}

\begin{question}
    A major benefit of securitization for a financial institution is the ability to remove assets from the balance sheet, which lowers risk and the required regulatory capital. While a large portion of the risk is removed from the balance sheet the originating financial institution often maintains a portion of the risk. Which of the following terms best identify the risk that is maintained by the originator?
\end{question}

\begin{choices}
    \item Correlation
    \item Excess spread
    \item First-loss piece
    \item Guarantor of collateral value
\end{choices}

\begin{correctanswer}
    C
\end{correctanswer}

\begin{explanation}
    \textbf{A选项错误。}
    这个说法不正确：
    \begin{itemize}
        \item 相关性不是发起人保留的风险
        \item 这是资产池的特征
        \item 这个选项与风险保留无关
    \end{itemize}

    \textbf{B选项错误。}
    这个说法不正确：
    \begin{itemize}
        \item 超额利差不是发起人保留的风险
        \item 这是信用增级方式
        \item 这个选项与风险保留无关
    \end{itemize}

    \textbf{C选项正确。}
    这个说法正确：
    \begin{itemize}
        \item 第一损失部分：
        \begin{itemize}
            \item 发起人通常保留所有权
            \item 信用质量最低的资产类别
            \item 违约时最先承担损失
            \item 代表发起人保留的风险
        \end{itemize}
    \end{itemize}

    \textbf{D选项错误。}
    这个说法不正确：
    \begin{itemize}
        \item 抵押品价值担保不是发起人保留的风险
        \item 这是外部信用增级方式
        \item 这个选项与风险保留无关
    \end{itemize}
\end{explanation}

\checklist{发起人通常保留第一损失部分的所有权}

\begin{question}
    Which of the following is an example of internal enhancements?
\end{question}

\begin{choices}
    \item CDS
    \item Put options on assets
    \item Overcollateralization
    \item Letters of credit
\end{choices}

\begin{correctanswer}
    C
\end{correctanswer}

\begin{explanation}
    \textbf{A选项错误。}
    这个说法不正确：
    \begin{itemize}
        \item CDS是外部信用增级
        \item 需要第三方参与
        \item 这个选项与内部增级无关
    \end{itemize}

    \textbf{B选项错误。}
    这个说法不正确：
    \begin{itemize}
        \item 资产看跌期权是外部信用增级
        \item 需要第三方参与
        \item 这个选项与内部增级无关
    \end{itemize}

    \textbf{C选项正确。}
    这个说法正确：
    \begin{itemize}
        \item 内部增级方式包括：
        \begin{itemize}
            \item 超额担保
            \item 直接股权发行
            \item 自留
            \item 现金担保账户
            \item 超额利差
        \end{itemize}
    \end{itemize}

    \textbf{D选项错误。}
    这个说法不正确：
    \begin{itemize}
        \item 信用证是外部信用增级
        \item 需要银行参与
        \item 这个选项与内部增级无关
    \end{itemize}
\end{explanation}

\checklist{超额担保是证券化产品的内部增级方式}

\begin{question}
    Under a 200\% PSA model, what is the single monthly mortality (SMM) for, respectively, month 10 and month 30?
\end{question}

\begin{choices}
    \item 0.12\% and 0.36\%
    \item 0.17\% and 0.51\%
    \item 0.28\% and 0.65\%
    \item 0.34\% and 1.06\%
\end{choices}

\begin{correctanswer}
    D
\end{correctanswer}

\begin{explanation}
    \textbf{D选项正确。}
    这个说法正确：
    \begin{itemize}
        \item 计算过程：
        \begin{itemize}
            \item 基础PSA模型：每月增加0.2\% CPR，直到第30个月达到6\%
            \item 第10个月：100\% PSA = 0.2\% $\times$ 10 = 2.0\%
            \item 第10个月：200\% PSA = 4.0\%
            \item SMM = 1-(1-4.0\%)$^{(1/12)}$ = 0.34\%
            \item 第30个月：100\% PSA = 0.2\% $\times$ 30 = 6.0\%
            \item 第30个月：200\% PSA = 12.0\%
            \item SMM = 1-(1-12.0\%)$^{(1/12)}$ = 1.06\%
        \end{itemize}
    \end{itemize}
\end{explanation}

\checklist{PSA模型下的月提前还款率}

\begin{question}
    The maximum benefits to the buyer of a credit-linked note (CLN) accrue when:
\end{question}

\begin{choices}
    \item There is a small credit downgrade
    \item There is no credit downgrade
    \item There is a large credit downgrade
    \item There is a default
\end{choices}

\begin{correctanswer}
    B
\end{correctanswer}

\begin{explanation}

    \textbf{B选项正确。}
    这个说法正确：
    \begin{itemize}
        \item CLN购买者的最大收益：
        \begin{itemize}
            \item 没有信用降级或违约时
            \item 获得高额回报
            \item 主要风险是降级或违约
            \item 违约可能导致血本无归
        \end{itemize}
    \end{itemize}


\end{explanation}

\checklist{信用联结票据在没有信用降级时获得最大收益}

\begin{question}
    Which of the following participants in the securitization process is least likely to face a conflict of interest?
\end{question}

\begin{choices}
    \item Credit rating agency and servicer
    \item Servicer and underwriter
    \item Custodian and trustee
    \item Trustee and manager
\end{choices}

\begin{correctanswer}
    C
\end{correctanswer}

\begin{explanation}
    \textbf{A选项错误。}
    这个说法不正确：
    \begin{itemize}
        \item 评级机构和服务商存在利益冲突
        \item 评级机构需要客观评级
        \item 服务商可能影响评级结果
    \end{itemize}

    \textbf{B选项错误。}
    这个说法不正确：
    \begin{itemize}
        \item 服务商和承销商存在利益冲突
        \item 服务商负责资产服务
        \item 承销商负责证券销售
    \end{itemize}

    \textbf{C选项正确。}
    这个说法正确：
    \begin{itemize}
        \item 托管人和受托人：
        \begin{itemize}
            \item 在证券化过程中作用最小
            \item 主要负责资产保管和监督
            \item 不参与核心业务决策
            \item 利益冲突最小
        \end{itemize}
    \end{itemize}

    \textbf{D选项错误。}
    这个说法不正确：
    \begin{itemize}
        \item 受托人和管理人存在利益冲突
        \item 受托人负责监督
        \item 管理人负责投资决策
    \end{itemize}
\end{explanation}

\checklist{托管人和受托人在证券化过程中利益冲突最小}


\begin{question}
    Which of the following measures are most likely to be used by a securitized product backed by student loans?
\end{question}

\begin{choices}
    \item Single monthly mortality (SMM), constant prepayment rate (CPR), and Public Securities Association (PSA)
    \item Loss curves and absolute prepayment speed (APS)
    \item Weighted average life (WAL), weighted average maturity (WAM), and weighted average coupon (WAC)
    \item Debt service coverage ratio (DSCR) and monthly payment rate (MPR)
\end{choices}

\begin{correctanswer}
    A
\end{correctanswer}

\begin{explanation}
    \textbf{A选项正确。}
    这个说法正确：
    \begin{itemize}
        \item 学生贷款证券化常用指标：
        \begin{itemize}
            \item 月提前还款率(SMM)
            \item 固定提前还款率(CPR)
            \item PSA提前还款模型
            \item 这些指标适用于学生贷款和抵押贷款
        \end{itemize}
    \end{itemize}

    \textbf{B选项错误。}
    这个说法不正确：
    \begin{itemize}
        \item 损失曲线和绝对提前还款速度
        \item 不适用于学生贷款证券化
        \item 这个选项与题目无关
    \end{itemize}

    \textbf{C选项错误。}
    这个说法不正确：
    \begin{itemize}
        \item 加权平均期限和加权平均票息
        \item 是资产池的基本特征
        \item 不是提前还款指标
    \end{itemize}

    \textbf{D选项错误。}
    这个说法不正确：
    \begin{itemize}
        \item 偿债覆盖率 and 月还款率
        \item 不适用于学生贷款证券化
        \item 这个选项与题目无关
    \end{itemize}
\end{explanation}

\checklist{学生贷款证券化使用SMM、CPR和PSA指标}


\begin{question}
    As the CRO of a retail bank, you are presenting to your Risk Committee the benefits of securitizing a pool of mortgages. Which of the following would you use to support your arguments that this will benefit the bank? 
    \begin{enumerate}
        \item It will improve your bank's return on capital
        \item Will immediately increase your bank's available capital
        \item You will be able to offer an attractive yield to investors
        \item It will lower your borrowing costs
    \end{enumerate}
\end{question}

\begin{choices}
    \item I and III only
    \item I, III, and IV only
    \item I, II, III and IV
    \item I, II, and IV only
\end{choices}

\begin{correctanswer}
    C
\end{correctanswer}

\begin{explanation}
    \textbf{A选项错误。}
    这个说法不正确：
    \begin{itemize}
        \item 只包含了部分好处
        \item 忽略了资本释放和融资成本降低
        \item 这个选项不完整
    \end{itemize}

    \textbf{B选项错误。}
    这个说法不正确：
    \begin{itemize}
        \item 只包含了部分好处
        \item 忽略了资本释放
        \item 这个选项不完整
    \end{itemize}

    \textbf{C选项正确。}
    这个说法正确：
    \begin{itemize}
        \item 证券化的所有好处：
        \begin{itemize}
            \item 提高资本回报率
            \item 立即增加可用资本
            \item 为投资者提供有吸引力的收益
            \item 降低融资成本
            \item 优化资产负债表结构
            \item 提高资产池透明度
        \end{itemize}
    \end{itemize}

    \textbf{D选项错误。}
    这个说法不正确：
    \begin{itemize}
        \item 只包含了部分好处
        \item 忽略了投资者收益
        \item 这个选项不完整
    \end{itemize}
\end{explanation}

\checklist{抵押贷款证券化可以带来资本回报、资本释放、投资者收益和成本降低}


\begin{question}
    Which of the following statements was not one of the flaws in the securitization process prior to the start of the credit crisis in 2007?
\end{question}

\begin{choices}
    \item An active originate-to-distribute model where a strong profit motive took precedence over ethical lending and underwriting
    \item The securitized products were so opaque that investors could not evaluate the true risks of the investment
    \item Structured investment vehicles (SIVs) were used to enhance the risk discovery process for investors and regulators
    \item Banks held securitized assets in off-balance-sheet entities, thus further masking the true risks in the system
\end{choices}

\begin{correctanswer}
    C
\end{correctanswer}

\begin{explanation}
    \textbf{A选项错误。}
    这个说法不正确：
    \begin{itemize}
        \item 这是2007年危机前的缺陷之一
        \item 利润动机优先于道德放贷
        \item 这个选项是真实缺陷
    \end{itemize}

    \textbf{B选项错误。}
    这个说法不正确：
    \begin{itemize}
        \item 这是2007年危机前的缺陷之一
        \item 产品不透明导致风险难以评估
        \item 这个选项是真实缺陷
    \end{itemize}

    \textbf{C选项正确。}
    这个说法正确：
    \begin{itemize}
        \item SIVs实际上：
        \begin{itemize}
            \item 增加了不透明度
            \item 是银行的表外实体
            \item 使投资者难以审查
            \item 不是用于增强风险发现
        \end{itemize}
    \end{itemize}

    \textbf{D选项错误。}
    这个说法不正确：
    \begin{itemize}
        \item 这是2007年危机前的缺陷之一
        \item 表外实体掩盖了真实风险
        \item 这个选项是真实缺陷
    \end{itemize}
\end{explanation}

\checklist{SIVs不是用于增强风险发现，而是增加了不透明度}


\begin{question}
    In a securitized transaction, over-collateralization results when:
\end{question}

\begin{choices}
    \item The originator puts aside some cash in a reserve account to absorb credit losses
    \item A securitization transaction carves up the cash flows generated from the asset pool into various pieces
    \item The interest payments and other fees received on the assets in the pool exceed the interest payment made on the asset-backed security (ABS) plus the fee paid to service the assets along with miscellaneous expenses
    \item The value of the assets in the pool exceeds the amount of asset-backed security(ABS) involved
\end{choices}

\begin{correctanswer}
    D
\end{correctanswer}

\begin{explanation}
    \textbf{A选项错误。}
    这个说法不正确：
    \begin{itemize}
        \item 这是现金储备账户的定义
        \item 用于吸收信用损失
        \item 不是超额抵押
    \end{itemize}

    \textbf{B选项错误。}
    这个说法不正确：
    \begin{itemize}
        \item 这是分级结构的定义
        \item 将现金流分成不同层级
        \item 不是超额抵押
    \end{itemize}

    \textbf{C选项错误。}
    这个说法不正确：
    \begin{itemize}
        \item 这是超额利差的定义
        \item 资产收益超过证券支出
        \item 不是超额抵押
    \end{itemize}

    \textbf{D选项正确。}
    这个说法正确：
    \begin{itemize}
        \item 超额抵押的定义：
        \begin{itemize}
            \item 资产池价值超过证券发行规模
            \item 实际发行的证券规模小于资产规模
            \item 是一种信用增级方式
        \end{itemize}
    \end{itemize}
\end{explanation}

\checklist{超额抵押是指资产池价值超过证券发行规模}



\begin{question}
    A pass-through mortgage-backed security (MBS) with a weighted average maturity (WAM) of 30 months has an original principle balance of \$2.0 billion. If the valuation model assumes a 300\% PSA prepayment speed, which is nearest to the first month's prepaid (not scheduled) principal?
\end{question}

\begin{choices}
    \item \$519,000
    \item \$833,000
    \item \$1.00 million
    \item \$2.15 million
\end{choices}

\begin{correctanswer}
    C
\end{correctanswer}

\begin{explanation}


    \textbf{C选项正确。}
    这个说法正确：
    \begin{itemize}
        \item 计算过程：
        \begin{itemize}
            \item 300\% PSA = 3 $\times$ CPR = 3 $\times$ 0.2\% = 0.6\%
            \item SMM = 1-(1-CPR)$^{(1/12)}$ = 1-(1-0.006)$^{(1/12)}$ = 0.00050138
            \item 提前还款金额 = 0.00050138 $\times$ 2,000,000,000 = 1,002,760
            \item 最接近100万美元
        \end{itemize}
    \end{itemize}


\end{explanation}

\checklist{300\%PSA下首月提前还款金额约为100万美元}


\begin{question}
    A risk consultant is reviewing a proposal for a US-based startup credit card company. A member of the senior management asks about the benefit of securitizing the company's loan book and the type of asset securitization structure that would best suit the credit card debt portfolio. If the consultant recommends securitization, which securitization structure would be most appropriate and why？
\end{question}

\begin{choices}
    \item An amortizing structure, because principal payments from the credit assets are paid to investors in a series of equal installments or in a lumpsum at maturities
    \item A master trust structure, because the structure allows multiple securitizations to be issued, and notes are issued out of the credit asset pool based on investor demand
    \item An amortizing structure, because the credit assets have relatively long maturities and low prepayment speeds
    \item A master trust structure, because the structure assumes no change in prepaymert speed, and investors are paid principal and interest on a coupon-by-coupon basis
\end{choices}

\begin{correctanswer}
    B
\end{correctanswer}

\begin{explanation}
    \textbf{A选项错误。}
    这个说法不正确：
    \begin{itemize}
        \item 信用卡债务不适合分期偿还结构
        \item 信用卡没有固定还款计划
        \item 这个选项与信用卡证券化无关
    \end{itemize}

    \textbf{B选项正确。}
    这个说法正确：
    \begin{itemize}
        \item 主信托结构的优势：
        \begin{itemize}
            \item 可以发行多个证券化产品
            \item 根据投资者需求发行票据
            \item 适合信用卡债务的特点
            \item 灵活性更高
        \end{itemize}
    \end{itemize}

    \textbf{C选项错误。}
    这个说法不正确：
    \begin{itemize}
        \item 信用卡债务期限不固定
        \item 提前还款速度不低
        \item 这个选项与信用卡证券化无关
    \end{itemize}

    \textbf{D选项错误。}
    这个说法不正确：
    \begin{itemize}
        \item 主信托结构不假设固定提前还款速度
        \item 不是按票息支付本息
        \item 这个选项描述不准确
    \end{itemize}
\end{explanation}

\checklist{信用卡债务证券化最适合使用主信托结构}


\begin{question}
    A fixed-income portfolio manager at a financial institution is estimating the weighted average coupon (WAC) of an MBS pool before making an investment. The pool consists of four residential mortgage loans. Information about the loans at the beginning of a specified month is given below:

    \begin{table}[H]
        \centering
        \begin{tabular}{|c|l|c|l|c|}
        \hline
        Loan & Outstanding beginning balance & Interest(coupon) rate & Scheduled principal payment & Remaining term (months) \\
        \hline
        1 & USD 240,000 & 6.5\% & USD 18,800 & 16 \\
        \hline
        2 & USD 525,000 & 5.0\% & USD 23,000 & 34 \\
        \hline
        3 & USD 470,000 & 6.8\% & USD 28,700 & 14 \\
        \hline
        4 & USD 875,000 & 3.0\% & USD 23,600 & 60 \\
        \hline
        \end{tabular}
    \end{table}

    What is the WAC of the MBS pool at the beginning of the month and what is the pool's ability to pay if the net coupon payable to investors in the MBS tranches is 4.76\%？
\end{question}

\begin{choices}
    \item WAC is 4.44\% and the pool's ability to pay is relatively weak
    \item WAC is 4.74\% and the pool's ability to pay is relatively weak
    \item WAC is 4.77\% and the pool's ability to pay is relatively strong
    \item WAC is 5.35\% and the pool's ability to pay is relatively strong
\end{choices}

\begin{correctanswer}
    B
\end{correctanswer}

\begin{explanation}
    \textbf{A选项错误。}
    这个说法不正确：
    \begin{itemize}
        \item WAC计算错误
        \item 低估了加权平均票息
        \item 这个选项完全错误
    \end{itemize}

    \textbf{B选项正确。}
    这个说法正确：
    \begin{itemize}
        \item 计算过程：
        \begin{itemize}
            \item 总余额 = 240,000 + 525,000 + 470,000 + 875,000 = 2,110,000
            \item WAC = (240,000×6.5\% + 525,000×5.0\% + 470,000×6.8\% + 875,000×3.0\%) / 2,110,000
            \item = 4.74\%
            \item 支付能力较弱，因为WAC(4.74\%) < 投资者票息(4.76\%)
        \end{itemize}
    \end{itemize}

    \textbf{C选项错误。}
    这个说法不正确：
    \begin{itemize}
        \item WAC计算错误
        \item 高估了支付能力
        \item 这个选项完全错误
    \end{itemize}

    \textbf{D选项错误。}
    这个说法不正确：
    \begin{itemize}
        \item WAC计算错误
        \item 高估了加权平均票息
        \item 这个选项完全错误
    \end{itemize}
\end{explanation}

\checklist{加权平均票息低于投资者票息时，资产池支付能力较弱}

\begin{question}
    A risk analyst at a hedge fund is assessing the following positions held by the fund:
    \begin{itemize}
        \item Contingent convertible bonds (CoCos) issued by GTU Bank; the triggering event for the conversion is when the price of GTU Bank's publicly traded stock reaches a specified threshold
        \item Senior tranche of an MBS issued by a bankruptcy-remote special purpose vehicle; the mortgages in the securitization are originated by GTU Bank
        \item Interest rate swap contract as the fixed-rate payer, with Bank NPV as the floating-rate payer; Bank NPV posts high-quality government bonds as collateral
    \end{itemize}
    The analyst considers a scenario under which interest rates increase and the stock price of GTU Bank decreases. As the analyst reviews the impact of these changes in market conditions on the three positions, which of the following is correct?
\end{question}

\begin{choices}
    \item The CoCos held by the hedge fund are less likely to be converted to stocks of GTU Bank
    \item The value of the senior tranche of the MBS held by the hedge fund will increase
    \item The hedge fund's wrong-way risk in the interest rate swap will decrease
    \item The value of the interest rate swap from the perspective of the hedge fund will increase
\end{choices}

\begin{correctanswer}
    D
\end{correctanswer}

\begin{explanation}
    \textbf{A选项错误。}
    这个说法不正确：
    \begin{itemize}
        \item 股价下跌会增加转换可能性
        \item 触发条件更容易满足
        \item 这个选项与实际情况相反
    \end{itemize}

    \textbf{B选项错误。}
    这个说法不正确：
    \begin{itemize}
        \item 利率上升会降低MBS价值
        \item 提前还款风险增加
        \item 这个选项与实际情况相反
    \end{itemize}

    \textbf{C选项错误。}
    这个说法不正确：
    \begin{itemize}
        \item 利率上升会增加错向风险
        \item 对手方信用风险增加
        \item 这个选项与实际情况相反
    \end{itemize}

    \textbf{D选项正确。}
    这个说法正确：
    \begin{itemize}
        \item 利率互换价值变化：
        \begin{itemize}
            \item 作为固定利率支付方
            \item 利率上升时价值增加
            \item 对手方提供高质量抵押品
            \item 信用风险较小
        \end{itemize}
    \end{itemize}
\end{explanation}

\checklist{利率上升时固定利率支付方的互换价值增加}




\newpage




\end{document}